% Part: second-order-logic
% Chapter: syntax-and-semantics
% Section: introduction

\documentclass[../../../include/open-logic-section]{subfiles}

\begin{document}

\olfileid[tr]{sol}{syn}{int}

\olsection{Giriş}

Birinci dereceden mantıkta, belirli bir dilin mantıksal olmayan simgelerini,
yani o dilin !!{constant}s, !!{function}s ve !!{predicate}s öğelerini,
birinci dereceden !!{structure}s hakkında şeyler ifade etmek üzere mantıksal
simgelerle birleştiririz. Bu, !!a{structure}~$\Struct{M}$ ile bir değişken
ataması~$s$'yi !!a{formula}~$!A$ ile ilişkilendiren sağlama kavramı
aracılığıyla yapılır: $\Sat{M}{!A}[s]$, ancak ve ancak $!A$'nın
!!{constant}s, !!{function}s ve !!{predicate}s öğeleri $\Struct{M}$'nin
belirttiği biçimde, serbest değişkenleri de~$s$'nin belirttiği biçimde
yorumlandığında ifade ettiği şey doğruysa geçerlidir. !!{identity}~$\eq$
yorumunun yanı sıra~$\lforall$ ve $\lexists$ yorumları da
$\Sat{M}{!A}[s]$ tanımına yerleştirilmiştir. İlki her zaman yapının
!!{domain}~$\Domain{M}$ üzerindeki özdeşlik bağıntısı olarak yorumlanır;
niceleyiciler de her zaman bütün !!{domain} üzerinde dolaşacak biçimde
yorumlanır. Fakat can alıcı nokta şudur: niceleme yalnızca !!{domain}
öğeleri üzerinde yapılabilir ve dolayısıyla bir niceleyicinin ardından
yalnızca nesne !!{variable}s gelebilir.

İkinci dereceden mantıkta hem dil hem de sağlama tanımı, serbest ve bağlı
fonksiyon ve yüklem değişkenlerini ve bunlar üzerindeki nicelemeyi içerecek
biçimde genişletilir. Bu değişkenlerin !!{function}s ve !!{predicate}s ile
ilişkisi, nesne değişkenlerinin !!{constant}s ile ilişkisiyle aynıdır.
İkinci dereceden mantığın terimlerinin ve !!{formula}s öğelerinin
oluşturulmasında aynı rolü oynarlar; üzerlerindeki niceleme de benzer biçimde
ele alınır. \emph{Standart} semantikte ikinci dereceden niceleyiciler, uygun
türdeki bütün olası nesneler üzerinde dolaşır (fonksiyon değişkenleri için
$\Domain{M}$'den~$\Domain{M}$'ye $n$-li fonksiyonlar, yüklem değişkenleri
için $n$-li bağıntılar). Örneğin,
$\lforall[\Obj{v_0}][(\Obj{P^1_0}(\Obj{v_0}) \lor \lnot
  \Obj{P^1_0}(\Obj{v_0}))]$ hem birinci hem de ikinci dereceden mantıkta bir
formülken, ikincisinde ayrıca
$\lforall[\Obj{V^1_0}][\lforall[\Obj{v_0}][(\Obj{V^1_0}(\Obj{v_0}) \lor
    \lnot \Obj{V^1_0}(\Obj{v_0}))]]$ ile
$\lexists[\Obj{V^1_0}][\lforall[\Obj{v_0}][(\Obj{V^1_0}(\Obj{v_0}) \lor
    \lnot \Obj{V^1_0}(\Obj{v_0}))]]$ ifadelerini de ele alabiliriz. Bunlar
hiçbir serbest değişken içermediğinden, ikinci dereceden mantığın
!!{sentence}s öğeleridir. Burada $\Obj{V^1_0}$, ikinci dereceden bir $1$-li
yüklem değişkenidir. $\Obj{V^1_0}$ için izin verilen yorumlar,
$\Obj{P^1_0}$ gibi $1$-li bir !!{predicate} öğesine verebileceğimiz yorumlarla,
yani~$\Domain{M}$'nin alt kümeleriyle aynıdır. O hâlde bunlar üzerindeki
niceleme,
$\lforall[\Obj{v_0}][(\Obj{V^1_0}(\Obj{v_0}) \lor \lnot
  \Obj{V^1_0}(\Obj{v_0}))]$ ifadesinin, $\Obj{V^1_0}$'a değer olarak
~$\Domain{M}$'nin bir alt kümesini atamanın bütün yollarında ya da en az bir
yolunda sağlandığını söylemek demektir. Her küme belirli bir nesneyi ya
içerdiğinden ya da içermediğinden, her iki tümce de her !!{structure}
içinde doğrudur.
\end{document}
